\section{Bounded delayed-screen comparison}
\label{app:delayed-screen}

We next test an established delayed-acceptance construction~\citep{banterle2019delayed}
while keeping the physical site-arc proposal and move-family schedule fixed.
For its complete MH log ratio $R$, choose a cheap bounded antisymmetric factor
$s(x,y,a)=-s(y,x,\bar a)$. The two acceptance probabilities are
\begin{equation}
 \alpha_1=\exp[\min(0,s)],\qquad
 \alpha_2=\exp[\min(0,R-s)].
\end{equation}
The potential is queried only after the first test passes. The product is no
larger than ordinary MH acceptance, but its forward/reverse ratio is $e^R$.
This standard correction preserves target invariance; it does not establish
mixing or the output law of a query-budget stopping rule. A zero factor exactly
recovers the previous random streams, queries and states. Geometry-valid
proposals and physically scored proposals are recorded separately.

The fixed cheap control uses squared covalent-radius-normalized bond stretch
with coefficient 10\,eV, nonbonded inverse-distance-power-12 repulsion with
coefficient 1\,eV, the declared COM restraint, and the known complete proposal
ratio. These define heuristic surrogate energies, not an additional physical
oracle. A linear screen fits four coefficients from zero; the neural screen
adds the difference of shared symmetric action-encoder readouts. An odd tanh
bounds all factors by $\log16$. Neither learned model loads earlier weights.
Two seeds each train for 300 steps with batch eight, using the same 36 fitting
parents, 12 internal-selection parents and all 1,671 attempts. No new physical
queries are used; inherited data construction still costs 18,110 raw calls.

The objective uses signed expected accepted potential decrease $J$ and expected
raw cost $C=\mathbb E[2\alpha_1]$, optimizing $J-\eta_0C$ with the original FIT
rate fixed. Both acceptance stages enter its gradient. Unsupported attempts
remain zero contributions in the full denominator. All final and control metrics
replay. Independent NumPy checks cover all 1,597 scored pairs for each of four
models, including both factors and their balance relation (maximum error
$3.09\times10^{-14}$). Six sensitive trained-head finite differences have maximum
error $4.93\times10^{-10}$.

\begin{table}[h]
\centering\small
\begin{tabular}{lrrr}
\toprule
Screen & $J$ (eV/attempt) & $C$ (calls/attempt) & $10^3 J/C$ \\
\midrule
No screen & 0.006182 & 1.9166 & 3.2256 \\
Fixed physical & 0.004513 & 0.7250 & 6.2251 \\
Linear, seed 0 & 0.005857 & 1.3861 & 4.2255 \\
Linear, seed 1 & 0.005896 & 1.2798 & 4.6073 \\
Neural, seed 0 & 0.005857 & 1.3861 & 4.2259 \\
Neural, seed 1 & 0.005896 & 1.2798 & 4.6073 \\
\bottomrule
\end{tabular}
\caption{Internal recorded-proposal estimates, retaining every failed attempt.
The last column is in $10^{-3}$\,eV per expected raw call. These are not achieved
query savings or measurements of screened chains.}
\label{tab:delayed-screen}
\end{table}

The learned screens improve the internal per-query proxy by about 31\% and
43\% over no screening; the fixed physical control's point improvement is 93\%.
Their learned-minus-physical intervals span zero and do not establish an
advantage. Neural-minus-linear differences are only $3.42\times10^{-7}$ and
$-4.81\times10^{-8}$\,eV per call, with intervals spanning zero. The additional
neural representation therefore has no established value in this comparison.
Intervals resample parents within four fixed compositions and are descriptive,
not multiplicity-adjusted. No complete-chain or wall-time benefit follows.
Expected rejected joint calls account for only 24.35\% of recorded FIT
complete-trajectory calls, before preparation and learning overhead; the large
conditional query-rate changes cannot be presented as a whole-sampler speedup.
The tested neural weights do not justify molecular scale-up.
